% VLDB WORKSHOP template version of 2024-05-15 enhances the ACM template, version 1.7.0:
% https://www.acm.org/publications/proceedings-template
% The ACM Latex guide provides further information about the ACM template.

\documentclass[sigconf, nonacm]{acmart}

\newcommand\vldbyear{2026}
\newcommand\vldbworkshop{Quality in Databases (QDB)}
\newcommand\vldbauthors{\authors}
\newcommand\vldbtitle{\shorttitle}
\newcommand\vldbavailabilityurl{}
\newcommand\vldbpagestyle{plain}

\begin{document}

\title[Medici\'on de credibilidad en redes sociales]{Medici\'on de credibilidad en plataformas de redes sociales mediante calidad de datos y provenance}

\author{Sebasti\'an Garc\'ia Parra}
\affiliation{%
  \institution{Universidad de la Rep\'ublica}
  \city{Montevideo}
  \country{Uruguay}
}
\email{TODO@example.org}

\begin{abstract}
Las plataformas de redes sociales se han convertido en fuentes primarias de
informaci\'on para decisiones de salud, pol\'itica, crisis y vida cotidiana. Sin
embargo, los usuarios suelen consumir y propagar contenido sin indicadores
objetivos sobre su credibilidad, especialmente cuando una publicaci\'on es copiada,
modificada o compartida entre plataformas. Este paper presenta un modelo de
calidad de datos para medir la credibilidad de clips en redes sociales. El modelo
organiza la credibilidad en dos dimensiones principales: trustworthiness, calculada
a partir de reputation, verifiability y expertise, y provenance, que captura el
origen y la evoluci\'on de un clip. Se introduce Trustworthiness Path Stability,
una m\'etrica de calidad que eval\'ua si el trustworthiness permanece estable a lo
largo del provenance. Tambi\'en se describe un flujo modular que normaliza clips
heterog\'eneos, genera features de dominio, calcula m\'etricas de calidad y
reconstruye provenance mediante interacciones expl\'icitas y equivalencia textual.
Una prueba de concepto sobre publicaciones acerca de estatinas y colesterol,
validada con un experto m\'edico, muestra alta concordancia entre el modelo y la
evaluaci\'on experta.
\end{abstract}

\maketitle

%%% do not modify the following VLDB block %%
%%% VLDB block start %%%
\pagestyle{\vldbpagestyle}
\begingroup\small\noindent\raggedright\textbf{VLDB Workshop Reference Format:}\\
\vldbauthors. \vldbtitle. VLDB \vldbyear\ Workshop: \vldbworkshop.\\
\endgroup
\begingroup
\renewcommand\thefootnote{}\footnote{\noindent
This work is licensed under the Creative Commons BY-NC-ND 4.0 International License. Visit \url{https://creativecommons.org/licenses/by-nc-nd/4.0/} to view a copy of this license. For any use beyond those covered by this license, obtain permission by emailing \href{mailto:info@vldb.org}{info@vldb.org}. Copyright is held by the owner/author(s). Publication rights licensed to the VLDB Endowment. \\
\raggedright Proceedings of the VLDB Endowment.
ISSN 2150-8097. \\
}\addtocounter{footnote}{-1}\endgroup
%%% VLDB block end %%%

%%% do not modify the following VLDB block %%
%%% VLDB block start %%%
\ifdefempty{\vldbavailabilityurl}{}{
\vspace{.3cm}
\begingroup\small\noindent\raggedright\textbf{VLDB Workshop Artifact Availability:}\\
The source code, data, and/or other artifacts have been made available at \url{\vldbavailabilityurl}.
\endgroup
}
%%% VLDB block end %%%

\section{Introducci\'on}

El contenido producido en redes sociales se utiliza cada vez m\'as como evidencia
para decisiones con consecuencias online y offline. Un usuario puede decidir si
comparte una publicaci\'on, conf\'ia en una fuente, acepta un consejo m\'edico o
rechaza un tratamiento a partir de informaci\'on cuya calidad es incierta. Este
riesgo aumenta por la velocidad con que las plataformas promueven la difusi\'on y
por la ausencia de est\'andares para preservar metadatos de calidad cuando el
contenido cruza fronteras entre plataformas.

Los mecanismos actuales, como cuentas verificadas o anotaciones comunitarias,
aportan se\~nales \'utiles pero no ofrecen una visi\'on general de calidad de datos
sobre la credibilidad. Adem\'as, los usuarios participan en varias plataformas y el
contenido puede ser copiado manualmente, citado, parafraseado o republicado sin
conservar los metadatos de la fuente original. As\'i, una misma pieza de
informaci\'on puede aparecer en varias versiones, con distintos autores, evidencias,
audiencias y contextos de confianza.

Este paper condensa una tesis de maestr\'ia sobre medici\'on de credibilidad en
plataformas de redes sociales. Su contribuci\'on principal es un modelo de calidad
de datos, sensible al provenance, para evaluar clips. Un clip representa una unidad
de informaci\'on social: post, tweet, comentario, caption de video o publicaci\'on
espec\'ifica de una plataforma. El modelo es agn\'ostico de la red social y adaptable
al dominio: un experto en calidad de datos define dimensiones y m\'etricas, mientras
que un experto de dominio ajusta ponderadores y umbrales.

El trabajo realiza tres contribuciones. Primero, define un modelo de calidad para
credibilidad que combina trustworthiness y provenance. Segundo, introduce
Trustworthiness Path Stability, una m\'etrica para capturar c\'omo cambia la
confianza cuando el contenido se propaga y transforma. Tercero, presenta un flujo
modular de procesamiento y una prueba de concepto en el dominio m\'edico, enfocada
en contenido sobre estatinas y colesterol.

\section{Antecedentes y Requisitos}

La credibilidad ha sido tratada como una calidad percibida y no como una propiedad
inherente de un objeto~\cite{FoggTseng1999}. Desde la calidad de datos, la
credibilidad se relaciona con trustworthiness, reputation y provenance
~\cite{BatiniScannapieco2016}. En redes sociales esta relaci\'on es especialmente
importante porque la evaluaci\'on ocurre bajo incertidumbre: el usuario ve solo una
parte del contexto original, y la autoridad aparente de una fuente puede resultar
enga\~nosa.

Trabajos sobre credibilidad en Twitter/X muestran que caracter\'isticas del mensaje,
del usuario, del tema y de la propagaci\'on ayudan a estimar credibilidad
~\cite{Castillo2011}. Otros trabajos abordan la calidad de datos en redes sociales
y adaptan dimensiones tradicionales de calidad a flujos de tweets
~\cite{Arolfo2020}. Estudios recientes muestran adem\'as que la expertise de la
fuente y las creencias previas influyen fuertemente en la evaluaci\'on de
publicaciones, en particular en dominios de salud~\cite{Kuutila2024}.

El provenance ofrece una mirada complementaria. El linaje de datos captura origen,
derivaci\'on y transformaciones, y ha sido reconocido como mecanismo central para
evaluar confianza en productos de datos~\cite{Simmhan2005}. PROV-SAID extiende
W3C PROV para representar difusi\'on de informaci\'on en redes sociales
~\cite{Taxidou2015}. Modela mensajes, agentes, influencia e interacciones, pero su
formulaci\'on original no modela directamente el cross-sharing ni m\'etricas de
calidad asociadas al camino de provenance.

De estas observaciones se derivan cuatro requisitos. El modelo debe ser agn\'ostico
de plataforma, porque las se\~nales var\'ian entre X, Facebook, TikTok, Reddit,
WhatsApp y otras redes. Debe ser sensible al provenance, porque la credibilidad
puede cambiar cuando el contenido se copia, revisa o comparte. Debe admitir
m\'etricas espec\'ificas de dominio, porque medicina, pol\'itica y emergencias
requieren evidencias distintas. Finalmente, debe ser modular para incorporar nuevos
generadores de features, m\'etricas y estrategias de b\'usqueda.

\section{Modelo de Calidad}

Modelamos \emph{Credibility} como un cluster compuesto por dos dimensiones:
\emph{Trustworthiness} y \emph{Provenance}. Trustworthiness estima la confiabilidad
esperada del clip dentro de su dominio. Provenance captura origen, derivaci\'on y
propagaci\'on.

Trustworthiness se descompone en tres factores. \emph{Reputation} eval\'ua la
percepci\'on externa del autor o fuente mediante m\'etricas como engagement y
confidence. \emph{Verifiability} eval\'ua si el clip aporta evidencia suficiente para
corroborar sus afirmaciones, incluyendo datos fuente, transparencia de
modificaciones y metadatos de verificaci\'on. \emph{Expertise} eval\'ua si el autor
tiene conocimiento espec\'ifico del dominio, usando m\'etricas como certification,
production, influence y experience.

Provenance se descompone en attribution, process y trustworthiness path.
Attribution identifica fuentes o entidades autoras. Process captura c\'omo se gener\'o
o modific\'o el clip. Trustworthiness path captura c\'omo evoluciona la confianza a lo
largo del camino de provenance.

\begin{table}[t]
  \caption{Dimensiones y factores principales del modelo de credibilidad.}
  \label{tab:model}
  \begin{tabular}{lll}
    \toprule
    Dimensi\'on & Factor & M\'etricas ejemplo\\
    \midrule
    Trustworthiness & Reputation & Engagement, confidence\\
    Trustworthiness & Verifiability & Access to source data\\
    Trustworthiness & Expertise & Certification, influence\\
    Provenance & Attribution & Attribution metadata\\
    Provenance & Process & Edit count, edit history\\
    Provenance & Path & Path stability\\
    \bottomrule
  \end{tabular}
\end{table}

Sea $T$ el valor de trustworthiness y $P$ el valor de provenance, ambos en
$[0,1]$. La credibilidad $C$ se calcula como una combinaci\'on ponderada:
\begin{equation}
  C = tT + pP,
\end{equation}
donde $t$ y $p$ son ponderadores definidos por dominio. Trustworthiness se calcula
de forma an\'aloga a partir de verifiability $V$, reputation $R$ y expertise $E$:
\begin{equation}
  T = vV + rR + eE.
\end{equation}
Los ponderadores pueden definirse manualmente por expertos, inferirse mediante
reglas o aprenderse a partir de clips etiquetados.

\section{Flujo con Provenance}

El flujo de procesamiento recibe un clip y calcula su credibilidad en seis etapas.
Primero, el clip se ingesta desde una fuente espec\'ifica. Segundo, se normaliza en
una representaci\'on JSON con identificadores, URI de plataforma, contenido, fechas,
geolocalizaci\'on, interacciones, multimedia, comentarios, m\'etricas de performance,
metadatos de usuario, features y m\'etricas de calidad. Tercero, los m\'odulos de
features extraen se\~nales espec\'ificas del dominio. Cuarto, los m\'odulos de calidad
calculan m\'etricas por factor. Quinto, el m\'odulo de provenance busca clips
equivalentes y reconstruye el grafo de linaje. Finalmente, el m\'odulo de
credibilidad combina las m\'etricas en un valor final.

El m\'odulo de provenance distingue interacciones expl\'icitas e impl\'icitas. Las
expl\'icitas son mecanismos propios de las plataformas: retweets, respuestas, citas o
compartidos. Las impl\'icitas ocurren cuando los usuarios copian, parafrasean o
republican contenido sin preservar un enlace directo. Estos casos son frecuentes en
cross-sharing, donde el contenido se mueve entre plataformas y puede perder sus
metadatos originales.

Para abordar las interacciones impl\'icitas, el flujo usa equivalencia entre clips.
Para clips de texto, instanciamos la equivalencia mediante similitud coseno con
TF-IDF~\cite{Yuntao2005}. Dados dos clips $c_1$ y $c_2$, su similitud textual es:
\begin{equation}
  Sim(c_1,c_2) = \frac{c_1 \cdot c_2}{\|c_1\|\|c_2\|}.
\end{equation}
Los clips equivalentes se representan luego mediante una estructura inspirada en
PROV-SAID. Adaptamos PROV-SAID sustituyendo mensajes por clips, agregando el
atributo \texttt{SocialNetwork}, modelando derivaciones entre plataformas y
asociando m\'etricas de calidad a las entidades del provenance.

\subsection{Trustworthiness Path Stability}

La principal m\'etrica derivada de provenance es \emph{Trustworthiness Path
Stability}. Mide si el trustworthiness permanece estable durante la propagaci\'on.
Sea $T_i$ el valor de trustworthiness del $i$-\'esimo clip en un camino de provenance
de largo $n$. Definimos:
\begin{equation}
  TPS = 1 - \frac{\sum_{i=1}^{n-1}|T_{i+1}-T_i|}{n-1}.
\end{equation}
Valores cercanos a 1 indican estabilidad. Valores bajos indican fluctuaciones,
posiblemente causadas por modificaciones del contenido, p\'erdida de calidad de la
fuente o propagaci\'on a trav\'es de actores menos confiables.

\section{Prueba de Concepto}

Se implement\'o una prueba de concepto para clips de salud relacionados con
estatinas y colesterol. El dominio fue elegido porque la desinformaci\'on sobre
estatinas puede influir en la adherencia a tratamientos y en decisiones de salud
p\'ublica. La implementaci\'on se despleg\'o en Google Cloud Platform con Pub/Sub
para ingesta, Cloud Storage para persistencia, Cloud Functions para m\'odulos de
features y calidad, y Cloud Composer para orquestaci\'on.

El experimento parti\'o de 97 clips recolectados en varias plataformas mediante
keywords sobre estatinas y colesterol. Se preseleccionaron once clips para cubrir
distintas redes y niveles de credibilidad. Un experto m\'edico clasific\'o los clips y
asign\'o puntajes.

Para el experimento, trustworthiness se calcul\'o con tres m\'etricas: access to source
data, certification y confidence. El experto asign\'o ponderadores de 0.4, 0.4 y 0.2,
respectivamente. Credibility combin\'o Trustworthiness Path Stability y
trustworthiness con ponderadores de 0.6 y 0.4 cuando hab\'ia provenance disponible.
Cuando no hab\'ia provenance, el modelo se apoy\'o solo en trustworthiness.

\begin{table}[t]
  \caption{Comparaci\'on entre juicio experto y salida del modelo.}
  \label{tab:results}
  \begin{tabular}{ccccc}
    \toprule
    Clip & Experto & Score & Modelo & Cred.\\
    \midrule
    1 & Medio & 0.60 & Alto & 0.80\\
    2 & Medio & 0.70 & Medio & 0.68\\
    3 & Medio & 0.60 & Medio & 0.68\\
    4 & Medio & 0.70 & Medio & 0.64\\
    5 & Medio & 0.60 & Alto & 0.78\\
    6 & Bajo & 0.00 & Bajo & 0.09\\
    7 & Bajo & 0.00 & Bajo & 0.18\\
    8 & Bajo & 0.00 & Bajo & 0.11\\
    9 & Bajo & 0.10 & Bajo & 0.09\\
    10 & Bajo & 0.00 & Bajo & 0.08\\
    11 & Bajo & 0.00 & Bajo & 0.36\\
    \bottomrule
  \end{tabular}
\end{table}

El modelo coincidi\'o con el nivel experto en nueve de los once clips
(\autoref{tab:results}). Las dos diferencias corresponden a clips donde el
provenance aument\'o la credibilidad: ambos ten\'ian caminos estables y hab\'ian sido
compartidos por usuarios con trustworthiness relativamente alto. El experto no
ten\'ia acceso a todos los perfiles involucrados en la propagaci\'on, por lo que su
puntaje se acerc\'o m\'as al trustworthiness aislado que al valor final sensible al
provenance.

La implementaci\'on preliminar obtuvo un tiempo promedio de ejecuci\'on de 5220 ms
por clip. Aunque el resultado es alentador, la muestra es peque\~na para realizar
afirmaciones sobre escalabilidad o precisi\'on predictiva. El valor del experimento
es mostrar que un modelo de calidad puede producir m\'etricas interpretables y
exponer cu\'ando el provenance cambia la evaluaci\'on.

\section{Discusi\'on}

El modelo propuesto desplaza la medici\'on de credibilidad desde una predicci\'on
monol\'itica hacia un flujo de calidad de datos. Esto tiene dos ventajas para QDB.
Primero, explicita los criterios: se puede inspeccionar si el score proviene de
verificabilidad, expertise, reputation o estabilidad del provenance. Segundo, separa
conocimiento de dominio e infraestructura. En medicina se puede usar certification
y access to source data, mientras que otros dominios pueden definir otras features y
umbrales.

Existen limitaciones. La prueba de concepto usa pocos clips y un \'unico experto de
dominio. La estrategia de equivalencia cubre solo texto y no incluye im\'agenes o
video. La reconstrucci\'on de provenance depende de APIs, buscadores y metadatos
p\'ublicos disponibles. Finalmente, el modelo asume que los expertos pueden definir
ponderadores significativos; trabajos futuros deber\'ian comparar calibraci\'on
manual, basada en reglas y aprendida.

Como trabajo futuro se propone escalar el experimento, incorporar equivalencia
multimodal, mejorar procesamiento en tiempo real, agregar metadatos de calidad
para completitud e incertidumbre del provenance, y dise\~nar visualizaciones de
m\'etricas para usuarios finales.

\section{Conclusi\'on}

Este paper present\'o un modelo de calidad de datos, sensible al provenance, para
medir credibilidad en plataformas de redes sociales. Al combinar trustworthiness,
provenance, equivalencia entre clips y Trustworthiness Path Stability, el modelo
ofrece una forma interpretable de evaluar contenido social entre plataformas y
dominios. La prueba de concepto sobre estatinas sugiere que el provenance puede
cambiar materialmente la evaluaci\'on de credibilidad y debe tratarse como una
se\~nal de calidad de primer orden.

\begin{acks}
TODO: Agregar financiamiento, tutor\'ia, agradecimientos institucionales y
agradecimientos al experto de dominio seg\'un corresponda.
\end{acks}

\bibliographystyle{ACM-Reference-Format}
\bibliography{references}

\end{document}
\endinput
